%% ============================================================
%%  Mechanistic Analysis of T2I Hallucinations via Mode Interpolation
%%  IEEE Conference Paper — Revised Version
%%  All changes from review integrated
%% ============================================================
\documentclass[conference]{IEEEtran}

\usepackage[T1]{fontenc}
\usepackage[utf8]{inputenc}
\usepackage{cite}
\usepackage{amsmath,amssymb,amsfonts}
\usepackage{graphicx}
\usepackage{textcomp}
\usepackage{xcolor}
\usepackage{booktabs}
\usepackage{multirow}
\usepackage{array}
\usepackage{url}
\usepackage{hyperref}
\usepackage{subcaption}
\usepackage{tabularx}

\hypersetup{
  colorlinks=true,
  linkcolor=blue!60!black,
  citecolor=green!50!black,
  urlcolor=blue!60!black
}

\newcommand{\T}[1]{\texttt{#1}}
\newcommand{\B}[1]{\textbf{#1}}

\def\BibTeX{{\rm B\kern-.05em{\sc i\kern-.025em b}\kern-.08em
    T\kern-.1667em\lower.7ex\hbox{E}\kern-.125emX}}

%% ============================================================
\begin{document}

\title{Mechanistic Analysis of Text-to-Image Hallucinations\\
via Mode Interpolation Under Classifier-Free Guidance}

\author{
  \IEEEauthorblockN{Hasini Jaishetty}
  \IEEEauthorblockA{
    231AI012\\
    \textit{Department of Information Technology}\\
    \textit{NITK Surathkal}\\
    \textit{jaishettyhasini.231ai012@nitk.edu.in}
  }
  \and
  \IEEEauthorblockN{Utkarsh Shukla}
  \IEEEauthorblockA{
    231AI040\\
    \textit{Department of Information Technology}\\
    \textit{NITK Surathkal}\\
    \textit{utkarshshukla.231ai040@nitk.edu.in}
  }
}

\maketitle

%% ============================================================
%%  ABSTRACT
%% ============================================================
\begin{abstract}
We investigate whether mode interpolation, a geometric mechanism
previously identified as the cause of hallucinations in unconditional
diffusion models~\cite{yi2024understanding}, persists as an explanatory
framework when semantic conditioning via cross-attention and
Classifier-Free Guidance (CFG) is applied in Text-to-Image (T2I)
generation. Using Stable Diffusion v1.4, we generate 18{,}000 images
across 220 distinct text prompts (8{,}000 internal and 10{,}000 external). Internally, the confirmed hallucination rate
collapses from 35.2\% at CFG\,=\,0 to exactly 0.0\% at CFG\,=\,7.5,
then partially rebounds to 1.3\% at CFG\,=\,15 which forms a
statistically confirmed U-curve (Spearman $\rho{=}{-}0.624$
descending, $p{\approx}0$; $\rho{=}{+}0.144$ ascending,
$p{=}2.23{\times}10^{-15}$). Trajectory variance achieves pooled
ROC-AUC\,=\,0.877 [95\% CI: 0.869,\,0.884]; critically, at CFG\,=\,1
trajectory variance achieves AUC\,=\,0.760 while CLIP score
collapses to AUC\,=\,0.522 (near-random), demonstrating that variance
discriminates hallucinations where CLIP cannot. A novel spatial
analysis of per-pixel 64$\times$64 variance heatmaps reveals a
three-tier taxonomy: hallucinated images exhibit uniform variance
spread (Spatial Concentration Index, SCI\,=\,0.386), geometric-only
failures exhibit localised hot-spots (SCI\,=\,0.486) and clean
images fall between (SCI\,=\,0.436). Cohen's $d{=}1.014$ for spatial
concentration is the largest effect size in our study. On 10{,}000
external images from 200 independent prompts using fixed internal
thresholds, trajectory variance AUC improves to 0.959 [95\% CI:
0.952,\,0.964] while CLIP AUC collapses to 0.475 [95\% CI:
0.419,\,0.533], indicating that CLIP is statistically indistinguishable from random chance on diverse prompts.
Non-overlapping bootstrap confidence intervals certify that trajectory
variance and CLIP score capture categorically different aspects of
generation quality.
\end{abstract}

\begin{IEEEkeywords}
diffusion models, hallucination detection, mode interpolation,
classifier-free guidance, trajectory variance, text-to-image
generation, spatial analysis, generalisation
\end{IEEEkeywords}

%% ============================================================
%%  1. INTRODUCTION
%% ============================================================
\section{Introduction}
\label{sec:introduction}

Denoising Diffusion Probabilistic Models (DDPMs) have become the
dominant paradigm for high-quality image synthesis~\cite{ho2020ddpm,
song2020ddim}. Despite impressive visual quality, these models
routinely produce hallucinated outputs with images that fall outside
the support of the true data distribution which manifest as duplicated
objects, physically impossible structures, anatomically incorrect
body parts and semantically incoherent scenes.

A landmark paper by Aithal et al.~\cite{yi2024understanding} provided
the first mechanistic geometric explanation for this phenomenon in
\emph{unconditional} diffusion models. The central claim is that
hallucinations arise from \emph{mode interpolation}: the learned
score function, being a smooth neural approximation rather than the
true score, generates samples in low-probability latent regions lying
geometrically \emph{between} valid data modes. That work proposed
\emph{trajectory variance} where the variance of predicted $\hat{x}_0$
estimates over the final reverse diffusion steps which acted as a practical
proxy for detecting these interpolating samples and achieving 95\%
detection accuracy on controlled distributions.

However, this theoretical framework was established exclusively for
unconditional generation. Modern T2I systems such as Stable
Diffusion~\cite{rombach2022ldm} fundamentally alter the generative
process through two mechanisms: (1)~semantic conditioning via
cross-attention between latent representations and CLIP-based text
embeddings~\cite{radford2021clip} and (2)~Classifier-Free
Guidance (CFG)~\cite{ho2022cfg} which amplifies the conditional
signal by interpolating between unconditional and conditional score
estimates with a user-controlled guidance scale $\lambda$. These
mechanisms may either suppress mode interpolation by constraining
trajectories to semantically valid regions or compound it by
introducing new forms of geometric instability when guidance is
too strong. A concurrent line of work investigates a mechanistically
distinct phenomenon, local generation bias in rendered text, which
is orthogonal to the latent-space geometric failures studied
here~\cite{liu2024texthallu}.

\textbf{Contributions.} This paper makes the following contributions:

\begin{enumerate}
  \item \textbf{CFG U-curve:} Empirical confirmation that trajectory
  variance forms a statistically significant U-curve as a function
  of CFG scale, with a minimum at CFG\,=\,7.5 (0.0\% hallucination
  rate) and a partial rebound at CFG\,=\,15 (1.3\%), while CLIP
  score increases monotonically throughout.

  \item \textbf{Per-CFG discriminability:} At CFG\,=\,1, trajectory
  variance achieves AUC\,=\,0.760 while CLIP collapses to
  AUC\,=\,0.522 (near-random), demonstrating that variance
  discriminates hallucinations precisely in the regime where CLIP
  provides no useful signal.

  \item \textbf{Three-tier spatial taxonomy:} Per-pixel variance
  heatmap analysis reveals that mode interpolation produces uniform
  spatial spread (SCI\,=\,0.386) while geometric-only failures
  produce localised hot-spots (SCI\,=\,0.486) which is a distinction
  invisible to the global scalar metric. Cohen's $d{=}1.014$ is
  the largest effect size in the study.

  \item \textbf{External generalisation:} Trajectory variance AUC
  improves from 0.877 to 0.959 on 200 unseen prompts with fixed
  thresholds; CLIP AUC collapses to 0.475 with a CI that includes
  0.5. Bootstrap confidence intervals confirm these signals are
  categorically different.

  \item \textbf{Dual hypothesis outcome:} Mode interpolation is
  CFG-dependent that is dominant at low guidance, completely suppressed
  at standard guidance, and partially reintroduced by
  over-guidance, while 24.3\% of all failures are attributable
  to semantic misalignment independent of geometric instability.
\end{enumerate}

%% ============================================================
%%  2. RELATED WORK
%% ============================================================
\section{Related Work}
\label{sec:related}

\subsection{Hallucination in Generative Models}
Hallucination in large language models has been extensively
studied~\cite{ji2023survey}. In vision-language models, hallucination
typically refers to generated content inconsistent with the visual
input~\cite{li2023evaluating}. For T2I models, the problem is less
formally defined: outputs may be semantically inconsistent with the
prompt, compositionally incorrect, or physically impossible.

\subsection{Geometric Mechanisms in Diffusion Models}
Aithal et al.~\cite{yi2024understanding} established that unconditional
diffusion model hallucinations arise from mode interpolation in the
learned score function, and proposed trajectory variance as a
computable proxy for geometric instability. Our work extends this
framework to conditioned T2I generation, testing whether the mechanism
persists under cross-attention and CFG, and introducing a spatial
decomposition of trajectory variance not considered in the original
work.

\subsection{Classifier-Free Guidance}
Ho and Salimans~\cite{ho2022cfg} introduced CFG as a method to trade
sample quality against diversity by amplifying the conditional signal
relative to the unconditional baseline. Our CFG U-curve finding
provides new evidence that there exists an optimal geometric stability
point (CFG\,=\,7.5) that is distinct from the maximum semantic
alignment point.

\subsection{Hallucination Detection}
Existing detection approaches for T2I hallucinations typically rely
on object detectors~\cite{wang2024yolov10} or vision-language
similarity metrics~\cite{radford2021clip} applied post-hoc. Our
work demonstrates that trajectory variance, computed during generation
at zero additional inference cost, generalises to unseen prompt
distributions where CLIP-based detection fails.

%% ============================================================
%%  3. METHODOLOGY
%% ============================================================
\section{Methodology}
\label{sec:methodology}

\subsection{Trajectory Variance Validation}
\label{sec:validation}

Before applying trajectory variance to Stable Diffusion, we first
validated the metric on four controlled toy distributions where
ground truth hallucinations are known exactly, following the
experimental methodology of~\cite{yi2024understanding}. The
validation criterion is a dead-zone-to-mode variance ratio exceeding
5.0 where failure would indicate the metric is unreliable.

\textbf{Setup.} Four experiments with increasing geometric complexity
(1D Gaussian mixtures with 3 modes and 5 modes, a 2D Gaussian grid and a simple shapes
generator) are used here with a fully-connected diffusion network. Trajectory variance
is computed over the final commitment-window steps of each denoising
trajectory.

\textbf{Results.} All four experiments passed: 3-Gaussian 1D
(ratio\,=\,16.14), 5-Gaussian 1D (ratio\,=\,8.80), 2D 3$\times$3
grid (ratio\,=\,5.55), Simple Shapes (ratio\,=\,17.43). All exceed
the threshold of 5.0, confirming that trajectory variance reliably
separates hallucinated from valid samples before application to
Stable Diffusion.

\subsection{Custom DDIM Pipeline}
\label{sec:pipeline}

The standard HuggingFace pipeline~\cite{wolf2020huggingface} does
not expose intermediate $\hat{x}_0$ estimates. We build a custom
DDIM loop~\cite{song2020ddim} for Stable Diffusion
v1.4~\cite{rombach2022ldm} that intercepts the per-step $\hat{x}_0$
prediction without altering the final generated image:
\begin{equation}
  \hat{x}_0(t) = \frac{x_t - \sqrt{1-\bar{\alpha}(t)}\,\hat{\varepsilon}_t}
                      {\sqrt{\bar{\alpha}(t)}}
  \label{eq:x0_pred}
\end{equation}
CFG is applied as a single UNet forward pass on the concatenated
$[\varnothing, c]$ batch:
\begin{equation}
  \hat{\varepsilon}_\lambda = \hat{\varepsilon}_\varnothing
    + \lambda\,(\hat{\varepsilon}_c - \hat{\varepsilon}_\varnothing)
  \label{eq:cfg}
\end{equation}
\textbf{Sanity check.} Standard and custom pipelines produce
pixel-identical outputs for the same seed (mean absolute pixel
difference\,=\,0.0109, which is 0.004\% of the full pixel range),
confirming the custom loop does not alter generation. See Figure~1.

\textbf{Variance window.} Trajectory variance is computed over the
last 15 DDIM steps, which we call the \emph{commitment window}
following~\cite{yi2024understanding}. This is the phase where the
model finalises its output structure. The cumulative noise schedule
at the start of this window is $\bar{\alpha}|_{\text{start}}{=}0.622$,
placing it in a valid measurement regime. The global scalar is:
\begin{equation}
  v = \text{Var}_t\bigl(\hat{x}_0(t)\bigr).\text{mean}()
  \label{eq:variance}
\end{equation}
For the spatial analysis (Section~\ref{sec:spatial}), we preserve
structure by computing channel-averaged 64$\times$64 variance maps
before the final spatial averaging.

\begin{figure}[t]
  \centering
  \includegraphics[width=\linewidth]{phase2_sanity_check.png}
  \caption{Pipeline sanity check. \textbf{Left:} Standard HuggingFace
  output. \textbf{Centre:} Custom DDIM loop output (same seed,
  var\,=\,0.002512). \textbf{Right:} Absolute pixel difference map
  (mean\,=\,0.0109, PASS).}
  \label{fig:sanity}
\end{figure}

\subsection{Experimental Dataset}
\label{sec:dataset}

\textbf{Prompt design.} Twenty prompts span five semantic categories
(Table~\ref{tab:prompts}), ranging from semantically precise
(Single-Concept: ``A single red apple'') to semantically impossible
(Conflicting: ``Frozen fire''). Each category targets a distinct
hallucination mechanism.

\begin{table}[h]
\centering
\caption{20-Prompt Design. 11 prompts validated by YOLO (AUTO);
9 prompts validated by CLIP only (MANUAL) because the required
concepts fall outside the COCO-80 vocabulary.}
\label{tab:prompts}
\small
\begin{tabular}{p{1.6cm}cp{3.0cm}c}
\toprule
\textbf{Category} & \textbf{n} & \textbf{Purpose} & \textbf{YOLO?} \\
\midrule
Single-Concept & 4 & Baseline; strict object count & Partial \\
Generic        & 4 & Underspecified; tests semantic breadth & Partial \\
Conflicting    & 4 & Impossible constraints; forces oscillation & No \\
Compositional  & 4 & Multi-object spatial binding & Yes \\
Negation       & 4 & Tests SD's ability to suppress attributes & Partial \\
\bottomrule
\end{tabular}
\end{table}

\textbf{Internal dataset.} There are two splits generated. The
\emph{calibration split} (1{,}000 images, seeds 50--99,
CFG\,=\,7.5 only) is used exclusively to compute per-category
detection thresholds and is never used in any AUC computation.
The \emph{experimental split} (7{,}000 images, seeds 0--49,
7 CFG values: $\{0, 1, 3, 5, 7.5, 10, 15\}$) is the main analysis
dataset. Critically, the same 50 seeds are reused across all 7 CFG
values, isolating the CFG effect from seed-level randomness.
All experiments use Stable Diffusion v1.4 with a DDIMScheduler
(50 steps) and they generated 512$\times$512 images. Trajectories are
stored as $(15, 4, 64, 64)$ floating-point arrays for each image.

\textbf{External dataset.} For generalisation testing, 10{,}000
images are generated from 200 independent external prompts (40 per
category) at CFG\,=\,7.5 only and with 50 seeds per prompt. All
detection thresholds from the internal calibration split are applied
\emph{unchanged} where no recalibration is performed. This fixed-threshold
constraint is what makes the external test a genuine test of
generalisation rather than re-calibration.

\subsection{Two-Signal Hallucination Labelling}
\label{sec:labelling}

\textbf{Per-category threshold calibration.} A global threshold would
systematically mislabel entire categories: Generic prompts have
naturally high variance (the model must choose between multiple valid
interpretations); Conflicting prompts have naturally low CLIP scores
because no image can satisfy contradictory constraints. We therefore
calibrate one threshold pair per category $k$ from the calibration
split:
\begin{align}
  \tau_\text{var}^{(k)} &= \mu_\text{var}^{(k)} + 2\sigma_\text{var}^{(k)}
  \label{eq:tvar} \\
  \tau_\text{clip}^{(k)} &= \mu_\text{clip}^{(k)} - 2\sigma_\text{clip}^{(k)}
  \label{eq:tclip}
\end{align}
Standard deviations use the population formula and at $n{=}200$ per
category the difference from the sample estimate is below 0.3\%.

\textbf{Labelling rule.} An image is labelled hallucinated only when
\emph{both} a geometric signal and a quality signal fail:
\begin{equation}
  \text{hallucinated} =
    \underbrace{\bigl(v > \tau_\text{var}^{(k)}\bigr)}_{\text{geometric}}
    \;\wedge\;
    \underbrace{\bigl(\text{yolo\_error} \vee \text{clip} < \tau_\text{clip}^{(k)}\bigr)}_{\text{quality}}
  \label{eq:label}
\end{equation}
The AND condition is essential: 364 images showed high variance but
passed quality checks (geometric-only failures), confirming that
geometric instability alone does not imply a wrong output.

This labelling produces four failure types: \textbf{both} (confirmed
hallucination), \textbf{geometric} (high variance, correct output),
\textbf{semantic} (normal variance, wrong output), and \textbf{none}
(clean generation).

\textbf{Validation tools.} YOLOv8x (confidence\,=\,0.25) is applied
to the 11 AUTO prompts with a strict count rule (exactly one object)
for Single-Concept prompts SC1 and SC2 only. CLIP ViT-L/14
(OpenCLIP, OpenAI pretrained) is applied to all 20 prompts and is
the sole signal for the 9 MANUAL prompts.

%% ============================================================
%%  4. RESULTS
%% ============================================================
\section{Results}
\label{sec:results}

\subsection{Hallucination Label Distribution}

Table~\ref{tab:distribution} summarises the 7{,}000 image labels.
62.3\% are clean, 8.2\% are confirmed hallucinations and 24.3\%
fail quality checks with \emph{normal} trajectory variance. This
last group is critical: it establishes semantic misalignment as
an independent failure mechanism and provides partial support for H$_0$.

\begin{table}[h]
\centering
\caption{Failure Type Distribution (7{,}000 Internal Images).}
\label{tab:distribution}
\small
\begin{tabular}{lrrl}
\toprule
\textbf{Type} & \textbf{n} & \textbf{\%} & \textbf{Interpretation} \\
\midrule
Clean (none)        & 4{,}364 & 62.3 & Passes all checks \\
Semantic only       & 1{,}698 & 24.3 & Stable traj.; wrong output \\
\textbf{Both (hallucinated)} & \textbf{574} & \textbf{8.2} & High var.\ + quality fail \\
Geometric only      &   364 &  5.2 & High var.; output correct \\
\bottomrule
\end{tabular}
\end{table}

\subsection{CFG Variance U-Curve}
\label{sec:ucurve}

Figure~\ref{fig:cfg_sweep} shows the central finding. Trajectory
variance is minimised at CFG\,=\,7.5 ($\bar{v}{=}0.002732$ and 0.0\%
hallucination rate) and rises in both directions. CLIP scores increase
monotonically with CFG ($\rho{=}{+}0.620$ and $p{\approx}0$) which confirms
that the two signals capture genuinely different aspects of generation
quality.

\begin{figure}[t]
  \centering
  \includegraphics[width=\linewidth]{figure2_cfg_sweep.png}
  \caption{CFG Sweep Analysis. \textbf{Left:} Mean trajectory variance
  ${\pm}$std per category vs.\ CFG. All categories show the U-shape.
  \textbf{Right:} Hallucination rate (left axis) and mean CLIP score
  (right axis). CLIP increases monotonically while hallucination rate
  forms a U.}
  \label{fig:cfg_sweep}
\end{figure}

\textbf{Mechanistic explanation.} At CFG\,=\,0, the null embedding
provides no directional guidance because the trajectories wander freely in latent
space. As CFG increases to 7.5, the conditional score provides
increasingly strong attraction toward the semantically aligned region which
suppresses mode-to-mode wandering. Above CFG\,=\,7.5, the guidance becomes
over-amplified: for semantically contradictory prompts
(e.g., ``Frozen fire'') as it forces the model to simultaneously satisfy
incompatible attractors, manifesting as renewed geometric oscillation.
The variance increase at high CFG is concentrated in the Conflicting
category (std\,=\,0.001717 at CFG\,=\,15 vs.\ 0.000286 for
Single-Concept). The ascending Spearman $\rho{=}{+}0.144$ is
statistically significant due to the large sample ($n{=}3{,}000$),
but the correlation magnitude is weak. We notice the rebound is real but modest.

\textbf{Internal consistency check.} At CFG\,=\,0, all 20 prompts
produce identical mean variance (0.003645) regardless of text content.
This confirms the null embedding is active and seed-level noise is the
sole variance source in that condition.

\subsection{Feature Space Analysis}

Figure~\ref{fig:scatter} shows all 7{,}000 images in the trajectory
variance $\times$ CLIP score feature space. The hallucinated cluster
(red) occupies the high-variance, low-CLIP quadrant as predicted by
Equation~(\ref{eq:label}). Semantic-only failures (blue) appear at
low CLIP across all variance levels, confirming stable trajectories
with wrong output. Geometric-only failures (orange) appear at high
variance but normal CLIP which confirms the geometric instability without
semantic failure.

\begin{figure}[t]
  \centering
  \includegraphics[width=\linewidth]{figure1_variance_vs_clip.png}
  \caption{Trajectory Variance vs.\ CLIP Score for all 7{,}000 images.
  Hallucinated images (red, $n{=}574$) cluster in the high-variance,
  low-CLIP quadrant. Geometric-only failures (orange, $n{=}364$)
  have high variance but normal CLIP.}
  \label{fig:scatter}
\end{figure}

\subsection{ROC-AUC and Per-CFG Analysis}
\label{sec:auc}

\textbf{Pooled AUC (Table~\ref{tab:roc}).} Both trajectory variance
(AUC\,=\,0.877) and CLIP (AUC\,=\,0.873) achieve good discrimination
on all 7{,}000 images pooled across CFG values. However, 548 of 574
hallucinated images (95.5\%) originate from CFG\,=\,0 and CFG\,=\,1.
The pooled AUC therefore reflects performance across the full guidance
range and should not be interpreted as production-setting performance.
The external evaluation (Section~\ref{sec:external}) is the primary
production-setting result.

\begin{table}[h]
\centering
\caption{ROC-AUC Results with Bootstrap 95\% CIs ($n{=}1{,}000$).
The combined detector uses logistic regression on variance and CLIP,
with 5-fold cross-validation (CV) to confirm stability across
different subsets of the data.}
\label{tab:roc}
\small
\begin{tabular}{lccc}
\toprule
\textbf{Signal} & \textbf{AUC} & \textbf{95\% CI} & \textbf{5-fold CV} \\
\midrule
Trajectory Variance    & 0.877 & [0.869, 0.884] & --- \\
CLIP Score             & 0.873 & [0.858, 0.888] & --- \\
Combined (LR)          & 0.898 & [0.886, 0.910] & $0.896{\pm}0.048$ \\
All Features (Var+SCI) & \textbf{0.926} & [0.918, 0.935] & $0.919{\pm}0.062$ \\
\midrule
Random baseline        & 0.500 & --- & --- \\
\bottomrule
\end{tabular}
\end{table}

\textbf{Per-CFG AUC (Table~\ref{tab:cfg_auc}).} Disaggregating AUC
by CFG level reveals the most important comparative finding of the
internal study. At CFG\,=\,0, both signals perform at similar levels
(variance 0.771, CLIP 0.761). At this setting, the null embedding
is active and hallucinated images are so dramatically wrong that even
a CLIP comparison (which should theoretically have no signal since
text conditioning is off) accidentally discriminates them due to
the severity of the failure. At CFG\,=\,1, where text conditioning
begins but is weak, trajectory variance maintains AUC\,=\,0.760
while CLIP collapses to AUC\,=\,0.522 (near-random). This is the
key finding of the per-CFG analysis: variance discriminates
hallucinations precisely in the regime where CLIP provides no
information. At CFG\,${\geq}$\,5, insufficient hallucinated images
exist for within-level AUC computation.

\begin{table}[h]
\centering
\caption{Per-CFG AUC. The pooled AUC\,=\,0.877 is driven by
CFG\,=\,0/1 (95.5\% of hallucinations).
$\star$\,=\,near-random; $\dagger$\,=\,small sample
($n_\text{hall}{<}15$), treat as indicative only.}
\label{tab:cfg_auc}
\small
\begin{tabular}{rrrcc}
\toprule
\textbf{CFG} & $\boldsymbol{n}_\text{hall}$ & $\boldsymbol{n}_\text{clean}$
& \textbf{Var AUC} & \textbf{CLIP AUC} \\
\midrule
0.0  & 352 & 648 & 0.771 & 0.761 \\
1.0  & 196 & 804 & 0.760 & $0.522^\star$ \\
3.0  &   9 & 991 & $0.906^\dagger$ & $0.696^\dagger$ \\
5.0  &   1 & 999 & --- & --- \\
7.5  &   0 & 1{,}000 & --- & --- \\
10.0 &   3 & 997 & $0.910^\dagger$ & $0.637^\dagger$ \\
15.0 &  13 & 987 & 0.849 & 0.517 \\
\bottomrule
\end{tabular}
\end{table}

\textbf{Per-category AUC.} Within individual prompt categories,
trajectory variance achieves strong discrimination: Negation\,=\,0.983,
Compositional\,=\,0.943, Single-Concept\,=\,0.924 (Figure~\ref{fig:roc}).
These values confirm that the metric reliably separates hallucinated
from clean images within a consistent prompt type.

\textbf{Statistical tests (Table~\ref{tab:stats}).}
The Mann-Whitney U test asks: across all possible pairs of one
hallucinated image and one clean image, does the hallucinated image
have higher variance? The answer is yes in 75.3\% of all pairs
($p{=}2.70{\times}10^{-197}$, effect size $r{=}0.753$).
Cohen's $d{=}0.819$ means the average hallucinated image has variance
0.82 standard deviations above the average clean image which is a large,
practically meaningful difference. Spearman correlations confirm
both arms of the U-curve are statistically significant. The
chi-squared test confirms that prompt category strongly predicts
failure type ($p{=}1.17{\times}10^{-120}$): Single-Concept and
Compositional prompts hallucinate far more than Negation, Generic,
and Conflicting.

\begin{table}[h]
\centering
\caption{Statistical Test Results.}
\label{tab:stats}
\small
\begin{tabular}{p{2.1cm}p{2.0cm}p{1.8cm}p{1.2cm}}
\toprule
\textbf{Test} & \textbf{Statistic} & \textbf{$p$-value} & \textbf{Effect} \\
\midrule
MW-U (Variance)       & $U{=}3.23{\times}10^6$ & $2.70{\times}10^{-197}$ & $r{=}0.753$ \\
MW-U (CLIP)           & $U{=}3.22{\times}10^6$ & $4.65{\times}10^{-193}$ & $r{=}0.745$ \\
Cohen's $d$ (Var)     & $d{=}0.819$            & ---                     & Large \\
Spearman (0$\to$7.5)  & $\rho{=}{-}0.624$      & ${\approx}0$            & Strong \\
Spearman (7.5$\to$15) & $\rho{=}{+}0.144$      & $2.23{\times}10^{-15}$  & Weak \\
Chi-sq.\ (category)   & $\chi^2{=}563.6$, $df{=}4$ & $1.17{\times}10^{-120}$ & $V{=}0.284$ \\
\bottomrule
\end{tabular}
\end{table}

\begin{figure}[t]
  \centering
  \includegraphics[width=\linewidth]{figure3_roc_curves.png}
  \caption{ROC Curve Analysis. \textbf{Left:} Overlaid ROC curves
  for trajectory variance (red), CLIP (blue), and combined logistic
  regression (green). \textbf{Right:} Per-category AUC for the three
  categories with sufficient hallucination count. Generic and
  Conflicting have zero confirmed hallucinations (see
  Section~\ref{sec:zero_rate}) and are excluded.}
  \label{fig:roc}
\end{figure}

%% ============================================================
%%  5. SPATIAL HALLUCINATION ANALYSIS
%% ============================================================
\section{Spatial Hallucination Analysis}
\label{sec:spatial}

\subsection{From Global Scalar to Spatial Heatmap}

All prior analysis used a single number per image, the global mean
variance. To understand \emph{where} the model was uncertain, not
just \emph{how much}, we compute per-pixel variance maps. For each
image, the $(15, 4, 64, 64)$ trajectory is reduced to a $64{\times}64$
variance map by computing variance across the 15 steps and averaging
over the 4 latent channels:
\begin{equation}
  M[x, y] = \frac{1}{4}\sum_{c=1}^{4}
             \text{Var}_{t}\bigl[\hat{x}_0(t)[c,x,y]\bigr]
  \label{eq:heatmap}
\end{equation}
For visualisation, this map is upsampled to $512{\times}512$ using
bilinear interpolation to match the output image resolution
(factor of 8, since $64{\times}8{=}512$). The statistics reported
in this section are all computed on the original $64{\times}64$ maps.
Consistency check: $\text{mean}(M)\,{=}\,v$ (the stored global scalar)
to within floating-point precision.

\textbf{Spatial metrics.} Two metrics characterise the
\emph{distribution} of instability rather than its magnitude:

The \emph{Spatial Concentration Index} (SCI) is the Gini coefficient
of pixel variance values:
\begin{equation}
  \text{SCI} = \frac{2\sum_{i=1}^{n} i\,m_{(i)}}{n\sum_{i=1}^{n} m_{(i)}}
               - \frac{n+1}{n}
  \label{eq:sci}
\end{equation}
SCI\,=\,0 means all pixels are equally unstable (uniform spread);
SCI\,=\,1 means all instability is in one pixel. Higher SCI indicates
more localised instability.

The \emph{Peak Variance Ratio} (PVR) measures hot-spot dominance:
\begin{equation}
  \text{PVR} = \frac{\text{mean}(m \geq P_{90})}{\text{mean}(m < P_{90})}
  \label{eq:pvr}
\end{equation}
Higher PVR indicates a dominant hot-spot; lower PVR indicates spread
instability.

\subsection{Three-Tier Spatial Taxonomy}
\label{sec:taxonomy}

Figure~\ref{fig:spatial_gallery} shows representative examples.
Table~\ref{tab:spatial} presents group statistics.

\begin{table}[h]
\centering
\caption{Spatial Metrics by Failure Type. Statistical tests use
Clean (halluc.\,=\,False, $n{=}6{,}426$); Clean (none, $n{=}4{,}364$)
is provided for reference. $\dagger$\,High PVR mean is driven by
Conflicting-category outliers at CFG\,${\geq}$\,10; median (4.58)
is more representative.}
\label{tab:spatial}
\small
\begin{tabular}{lrrrr}
\toprule
\textbf{Group} & \textbf{n} & \textbf{SCI} & \textbf{PVR mean} & \textbf{PVR med.} \\
\midrule
Hallucinated (both) & 574   & 0.386 & 3.33 & 3.25 \\
Semantic only       & 1{,}698 & 0.410 & 3.65 & 3.53 \\
Geometric only      & 364   & \textbf{0.486} & $5.82^\dagger$ & $4.58^\dagger$ \\
Clean (none)        & 4{,}364 & 0.442 & 4.23 & 3.95 \\
Clean (hall.\,=\,F) & 6{,}426 & 0.436 & 4.17 & --- \\
\bottomrule
\end{tabular}
\end{table}

\textbf{Key finding.} Geometric-only failures have the \emph{highest}
SCI meaning the most localised instability, which was the opposite of our
initial hypothesis. This reveals two mechanistically distinct types
of high-variance images that are indistinguishable from the global
scalar alone. Full hallucination (mode interpolation between two
entirely different global images) produces \emph{uniform} instability:
because the two candidate images differ everywhere, every pixel is
equally uncertain. Geometric-only failures (where the model commits
globally but cannot resolve one irresolvable semantic region) produce
a \emph{focused hot-spot} while the rest of the image is stable.

\textbf{Reading a heatmap.} A heatmap alone does not determine failure
type as it must be combined with the global variance scalar. Images with
high global variance \emph{and} low SCI (uniform purple haze, no
dominant bright region) are confirmed hallucinations. Images with high
global variance \emph{and} high SCI (clear bright dot against near-black
background) are geometric-only failures. Clean images exhibit low global
variance where any hot-spots present correspond to fine structural decisions
(object contours, reflections) and are expected, not pathological.
This three-way combination: magnitude, spread, and location is what
distinguishes the failure types.

\begin{figure}[t]
  \centering
  \includegraphics[width=\linewidth]{figure5_spatial_gallery.png}
  \caption{Spatial variance heatmaps by failure type.
  \textbf{Top row (Hallucinated):} Uniform purple haze---global mode
  interpolation; every pixel is uncertain.
  \textbf{Middle row (Geometric-only):} Localised bright hot-spot;
  global structure committed but one region unresolvable.
  \textbf{Bottom row (Clean):} Moderate diffuse pattern following
  the object boundary.}
  \label{fig:spatial_gallery}
\end{figure}

\textbf{Statistical significance (Table~\ref{tab:spatial_stats}).}
Cohen's $d{=}{-}1.014$ for SCI is the largest effect size in the
entire study, exceeding the global variance result ($d{=}0.819$).

\begin{table}[h]
\centering
\caption{Spatial Statistical Tests. Hallucinated ($n{=}574$) vs
clean/halluc.\,=\,False ($n{=}6{,}426$).}
\label{tab:spatial_stats}
\small
\begin{tabular}{p{1.8cm}p{1.0cm}p{1.0cm}p{3.2cm}}
\toprule
\textbf{Test} & \textbf{Hall.\ mean} & \textbf{Clean mean} & \textbf{$p$-value / Effect} \\
\midrule
MW-U (SCI)        & 0.386 & 0.436 & $8.08{\times}10^{-147}$; $r{=}0.308$ \\
MW-U (PVR)        & 3.333 & 4.165 & $2.06{\times}10^{-148}$; $r{=}0.310$ \\
Cohen's $d$ (SCI) & ---   & ---   & $d{=}{-}1.014$ (\textbf{largest}) \\
Cohen's $d$ (PVR) & ---   & ---   & $d{=}{-}0.763$ (Large) \\
\bottomrule
\end{tabular}
\end{table}

Figure~\ref{fig:sci_dist} shows the SCI distributions for all four
groups.

\begin{figure}[t]
  \centering
  \includegraphics[width=\linewidth]{figure6_sci_distributions.png}
  \caption{SCI Distribution by Failure Type. \textbf{Left:} Density
  plot. Hallucinated (red) peaks at lower SCI than clean (green).
  Geometric-only (orange) is the most right-shifted.
  \textbf{Right:} Per-category comparison.
  All differences $p{<}0.001$ ($***$).}
  \label{fig:sci_dist}
\end{figure}

\subsection{AUC with Spatial Features}

SCI and PVR alone achieve AUC\,=\,0.824 and 0.826 respectively (this is
slightly below global variance (0.877) ) because they measure
instability \emph{shape} rather than magnitude. Combined with global
variance and CLIP, the all-features logistic regression reaches
AUC\,=\,0.926 [CI: 0.918, 0.935], confirmed stable by 5-fold
cross-validation (0.919\,${\pm}$\,0.062).

%% ============================================================
%%  6. EXTERNAL VALIDATION
%% ============================================================
\section{External Validation}
\label{sec:external}

\subsection{Design: Fixed-Threshold Constraint}

The scientific value of external validation depends on one constraint which are
the thresholds calibrated from the internal dataset
($\tau_\text{var}^{(k)}$, $\tau_\text{clip}^{(k)}$) are applied to
all external images \emph{without recalibration}. If thresholds were
re-fitted to the external data then any AUC improvement would reflect
recalibration rather than genuine generalisation. This design mirrors
real deployment where a threshold is set once on a reference
distribution and then applied continuously to new user-generated
content.

\textbf{Calibration verification.} At CFG\,=\,7.5 we notice that external mean
variance (0.002754) differs from the internal mean (0.002732) by
only 0.8\%, and external mean CLIP (0.2406) differs from internal
(0.2401) by 0.2\%. Any observed AUC change therefore reflects genuine
generalisation and not the calibration error.

\subsection{Generalisation Results}

\begin{table}[h]
\centering
\caption{Internal vs.\ External AUC. CLIP CI [0.419, 0.533]
includes 0.50: CLIP is not statistically distinguishable from
random chance on diverse prompts. Variance CI [0.952, 0.964]
excludes 0.50 entirely.}
\label{tab:gen}
\small
\begin{tabular}{lcccc}
\toprule
\textbf{Signal} & \textbf{Int.\ AUC} & \textbf{Ext.\ AUC}
& \textbf{Ext.\ CI} & \textbf{Result} \\
\midrule
Variance & 0.877 & \textbf{0.959} & [0.952, 0.964] & \checkmark \\
CLIP     & 0.873 & 0.475 & [0.419, 0.533] & $\times^\star$ \\
Combined & 0.898 & \textbf{0.960} & [0.951, 0.967] & \checkmark \\
\bottomrule
\multicolumn{5}{l}{$\star$\,CI includes 0.50: not distinguishable from random chance.}
\end{tabular}
\end{table}

Figure~\ref{fig:gen} shows the full comparison. Three findings stand
out.

\textbf{Variance generalises and improves.} AUC increases from 0.877
to 0.959. Cohen's $d$ improves from 0.82 to 1.879. A big insight from this is that now the effect size has
nearly doubled. Hallucinated images on external data show slightly
higher mean variance (0.003994 vs.\ 0.003645), consistent with the
external prompts being harder to satisfy semantically.

\textbf{CLIP fails categorically.} External CLIP AUC\,=\,0.475
with CI\,=\,[0.419, 0.533] which includes 0.50. CLIP is not
statistically distinguishable from random chance on diverse prompts.
The reason is structural: abstract Generic prompts (``A warm place'',
``Something round'') and Conflicting prompts (``Dry ocean'',
``Cold fire'') produce intrinsically low CLIP scores for any generated
image regardless of quality. The CLIP distributions for hallucinated
and clean images overlap almost completely on these categories and
destroying the discriminative power.

\textbf{CI non-overlap.} Variance external CI [0.952, 0.964] does
not overlap with CLIP external CI [0.419, 0.533] by a margin of
0.419. This non-overlap is the definitive statistical statement that
these two signals are drawn from categorically different AUC
distributions on diverse prompt data.

\textbf{Why variance generalises and CLIP does not.} Trajectory
variance measures a property of the diffusion process itself that is how
much the model's prediction of the final image changed during the
last 15 denoising steps. This computation is independent of what
text was given. A model oscillating between two globally incompatible
image structures will produce high variance whether the prompt is
``A red apple'' or ``The sound of emptiness.''

But the CLIP cosine similarity, by contrast, is calibrated against specific
text embeddings. The internal thresholds were set for 20 specific
prompts. External dataset used abstract and impossible prompts resulting in inherently low CLIP
scores for any generated image and making the hallucinated/clean
threshold meaningless on diverse prompts. Trajectory variance requires
only a model-level calibration thus establishing the normal variance
range for this specific diffusion model and then it generalizes across
semantic domains.

\begin{figure}[t]
  \centering
  \includegraphics[width=\linewidth]{figure9_generalisation.png}
  \caption{External Validation Results. \textbf{Top-left:}
  Variance improves ($0.877{\to}0.959$); CLIP collapses
  ($0.873{\to}0.475$). \textbf{Top-centre:} ROC curves.
  \textbf{Top-right:} Variance distributions; Cohen's $d{=}1.879$.
  \textbf{Bottom-left:} Per-category AUC. Variance ($>$0.95)
  across all testable categories; CLIP fails on abstract prompts.}
  \label{fig:gen}
\end{figure}

%% ============================================================
%%  7. BOOTSTRAP CONFIDENCE INTERVALS
%% ============================================================
\section{Bootstrap Confidence Intervals}
\label{sec:bootstrap}

The AUC estimates are accompanied by 1{,}000-iteration percentile
bootstrap 95\% confidence intervals (seed\,=\,42, all
$n_{\text{valid}}{=}1{,}000$). Combined models refit the logistic
regression on each bootstrap resample.

Table~\ref{tab:ci} presents all CIs. The CI widths are informative:
variance external has the narrowest CI (width\,=\,0.012) because
high AUC with consistent class separation is stable across
resamples. CLIP external has the widest CI (width\,=\,0.114)
because AUC near 0.5 with a noisy signal varies widely across
resamples. Here, all-features CI [0.918, 0.935] does not overlap
with the combined (Var+CLIP) CI lower bound (0.886), confirming that
adding spatial features provides a statistically meaningful AUC
improvement.

\begin{table}[h]
\centering
\caption{All Bootstrap 95\% CIs ($n{=}1{,}000$, seed\,=\,42).}
\label{tab:ci}
\small
\begin{tabular}{lcccc}
\toprule
\textbf{Signal} & \textbf{AUC} & \textbf{Lower} & \textbf{Upper} & \textbf{Width} \\
\midrule
Var Internal       & 0.877 & 0.869 & 0.884 & 0.016 \\
CLIP Internal      & 0.873 & 0.858 & 0.888 & 0.030 \\
Combined (Var+CLIP)& 0.898 & 0.886 & 0.910 & 0.024 \\
SCI alone          & 0.824 & 0.809 & 0.841 & 0.032 \\
PVR alone          & 0.826 & 0.811 & 0.842 & 0.031 \\
All Features       & 0.926 & 0.918 & 0.935 & 0.018 \\
\midrule
\textbf{Var External}  & \textbf{0.959} & \textbf{0.952} & \textbf{0.964} & 0.012 \\
\textbf{CLIP External} & \textbf{0.475} & \textbf{0.419} & \textbf{0.533} & 0.114 \\
Combined External  & 0.960 & 0.951 & 0.967 & 0.015 \\
\bottomrule
\end{tabular}
\end{table}

Figure~\ref{fig:bootstrap} shows the bootstrap CI bar charts.

\begin{figure}[t]
  \centering
  \includegraphics[width=\linewidth]{figure10_bootstrap_ci.png}
  \caption{Bootstrap 95\% CIs on all AUC estimates.
  \textbf{Left:} Internal signals.
  \textbf{Right:} Internal vs.\ external. Variance pair
  (blue/teal) overlaps---consistent generalisation. CLIP pair
  (purple/red) diverges dramatically---CLIP is
  distribution-specific.}
  \label{fig:bootstrap}
\end{figure}

%% ============================================================
%%  8. DISCUSSION
%% ============================================================
\section{Discussion}
\label{sec:discussion}

\subsection{Hypothesis Assessment}

\textbf{H$_1$ (mode interpolation persists): SUPPORTED with
qualification.} At CFG\,=\,0--1 we observe the hallucination rates of 35.2\% and
19.6\% which confirms that text conditioning alone does not prevent mode
interpolation when guidance is absent or weak. Trajectory variance
achieves AUC\,=\,0.760 at CFG\,=\,1 while CLIP is at 0.522.
External AUC\,=\,0.959 with Cohen's $d{=}1.879$ confirms robust
generalisation. The qualification: at CFG\,${\geq}$\,3, hallucination
rate drops below 1\%, and at CFG\,=\,7.5 it is exactly 0.0\%.
H$_1$ is most strongly supported at low CFG.

\textbf{H$_0$ (semantic misalignment dominates): PARTIALLY SUPPORTED.}
1{,}698 images (24.3\%) fail quality checks with normal trajectory
variance. Negation violations (46.9\% YOLO error for ``Face, no eyes'')
are driven by Stable Diffusion ignoring text constraints and not the geometric
instability. The CLIP AUC collapse on external data (CI includes 0.5)
confirms that semantic failure mechanisms are prompt-distribution
specific.

\subsection{The Zero-Hallucination Categories}
\label{sec:zero_rate}

Generic and Conflicting prompts both show 0\% confirmed hallucination
rate. This part requires explanation. These 2{,}800 images (40\% of the
dataset) contribute to the clean pool in AUC computation, making the
overall AUC harder to achieve meaning 0.877 is a conservative estimate.

For \textbf{Generic} prompts, the model resolves underspecification
by cleanly committing to one interpretation (generating a single fruit
rather than oscillating between multiple). The per-category threshold
$\tau_\text{var}^\text{Gen}{=}0.005510$ is the highest of all
categories which is calibrated from a naturally noisier baseline. The 13
geometric-only cases indicate some instability but this is not enough to
cross the calibrated threshold. This is a correct calibrated outcome.

For \textbf{Conflicting} prompts (all MANUAL, CLIP-only validation),
the threshold $\tau_\text{var}^\text{Con}{=}0.005774$ is calibrated
from a baseline where impossible concepts already produce noisier
trajectories. While 302 Conflicting images fail CLIP quality checks
(semantic failure), they rarely exceed the category-specific variance
threshold. However the spatial analysis reveals that geometric-only
Conflicting images at CFG\,${\geq}$\,10 produce extreme localised
hot-spots (mean PVR\,=\,12.64 for this subgroup) which is the highest of any
subgroup in the study and it is confirming that this category does produce
geometric instability, but its spatially localised nature prevents it
from crossing the global variance threshold.

\subsection{Implications for T2I System Design}

\textbf{CFG selection.} CFG\,=\,7.5 is the geometric stability
optimum. Below CFG\,=\,3, mode interpolation is a significant concern.
Above CFG\,=\,10, semantically complex or contradictory prompts exhibit
renewed instability.

\textbf{Runtime monitoring.} Trajectory variance can be monitored as
an unsupervised quality signal at zero additional inference cost during
normal DDIM generation. Images with variance exceeding 2$\sigma$ above
the category-level mean warrant review.

\textbf{Detector deployment.} CLIP-based detection requires
prompt-specific calibration and fails on semantically diverse or
abstract prompts. For production systems which handles varied user inputs, we conclude that the
trajectory variance is the superior signal. The combined detector
(variance\,+\,SCI) achieves AUC\,=\,0.926 internally with stable
cross-validation and is the recommended configuration.

%% ============================================================
%%  9. CONCLUSION
%% ============================================================
\section{Conclusion}
\label{sec:conclusion}

We have presented an empirical investigation of mode interpolation
as an explanatory mechanism for T2I hallucinations, spanning 18{,}000
generated images across 220 distinct prompts.

\textbf{At the generation level:} mode interpolation is real and
measurable in T2I but is strongly CFG-dependent. At CFG\,=\,7.5,
conditioning completely suppresses mode interpolation (0.0\% rate).
Over-guidance reintroduces a new instability for semantically
contradictory prompts. The 24.3\% semantic-only failure rate confirms
that text encoder limitations are an independent failure mechanism.

\textbf{At the spatial level:} per-pixel variance heatmaps reveal
a three-tier taxonomy: uniform spread (hallucinated, SCI\,=\,0.386),
localised hot-spot (geometric-only, SCI\,=\,0.486) and diffuse
moderate (clean, SCI\,=\,0.436). Cohen's $d{=}1.014$ for spatial
concentration is the largest effect size in the study. This taxonomy
is invisible to the global scalar metric.

\textbf{At the generalisation level:} trajectory variance AUC
improves from 0.877 to 0.959 on 200 unseen prompts with fixed
thresholds; Cohen's $d$ nearly doubles from 0.82 to 1.879. CLIP
AUC collapses from 0.873 to 0.475 with a CI that includes 0.5.
Bootstrap CI non-overlap statistically certifies that trajectory
variance is a distribution-agnostic geometric detector while CLIP
is a prompt-specific semantic proxy. This
distinction has direct practical consequences in the T2I deployment

%% ============================================================
%%  REFERENCES
%% ============================================================
\bibliographystyle{IEEEtran}

\begin{thebibliography}{15}\bibitem{yi2024understanding}S.~K.~Aithal, P.~Maini, Z.~C.~Lipton, and J.~Z.~Kolter,``Understanding hallucinations in diffusion models through modeinterpolation,''in \textit{Advances in Neural Information Processing Systems (NeurIPS)}, 2024.

\bibitem{ho2020ddpm}J.~Ho, A.~Jain, and P.~Abbeel,``Denoising diffusion probabilistic models,''in \textit{Advances in Neural Information Processing Systems (NeurIPS)},vol.~33, pp.~6840--6851, 2020.

\bibitem{song2020ddim}J.~Song, C.~Meng, and S.~Ermon,``Denoising diffusion implicit models,''in \textit{International Conference on Learning Representations (ICLR)},2021.

\bibitem{rombach2022ldm}R.~Rombach, A.~Blattmann, D.~Lorenz, P.~Esser, and B.~Ommer,``High-resolution image synthesis with latent diffusion models,''in \textit{Proceedings of the IEEE/CVF CVPR}, pp.~10684--10695, 2022.

\bibitem{radford2021clip}A.~Radford, J.~W.~Kim, C.~Hallacy, A.~Ramesh, G.~Goh, S.~Agarwal,G.~Sastry, A.~Askell, P.~Mishkin, J.~Clark, G.~Krueger, andI.~Sutskever,``Learning transferable visual models from natural language supervision,''in \textit{Proceedings of ICML}, pp.~8748--8763, 2021.

\bibitem{ho2022cfg}J.~Ho and T.~Salimans,``Classifier-free diffusion guidance,''in \textit{NeurIPS Workshop on Deep Generative Models and DownstreamApplications}, 2022.

\bibitem{liu2024texthallu}R.~Lu, R.~Wang, K.~Lyu, X.~Jiang, G.~Huang, and M.~Wang,``Towards understanding text hallucination of diffusion models vialocal generation bias,''in \textit{International Conference on Learning Representations (ICLR)}, 2025.

\bibitem{ji2023survey}Z.~Ji, N.~Lee, R.~Frieske, T.~Yu, D.~Su, Y.~Xu, E.~Ishii, Y.~Bang,A.~Madotto, and P.~Fung,``Survey of hallucination in natural language generation,''\textit{ACM Computing Surveys}, vol.~55, no.~12, pp.~1--38, 2023.

\bibitem{li2023evaluating}Y.~Li, Y.~Du, K.~Zhou, J.~Wang, W.~X.~Zhao, and J.-R.~Wen,``Evaluating object hallucination in large vision-language models,''in \textit{Proceedings of EMNLP}, 2023.

\bibitem{wang2024yolov10}A.~Wang, H.~Chen, L.~Liu, K.~Chen, Z.~Lin, J.~Han, and G.~Ding,``YOLOv10: Real-time end-to-end object detection,''\textit{arXiv preprint arXiv:2405.14458}, 2024.

\bibitem{wolf2020huggingface}T.~Wolf, L.~Debut, V.~Sanh, J.~Chaumond, C.~Delangue, A.~Moi,P.~Cistac, T.~Rault, R.~Louf, M.~Funtowicz, et~al.,``HuggingFace's Transformers: State-of-the-art natural languageprocessing,''in \textit{Proceedings of EMNLP: System Demonstrations}, pp.~38--45,2020.

\bibitem{saharia2022photorealistic}C.~Saharia, W.~Chan, S.~Saxena, L.~Li, J.~Whang, E.~L.~Denton,K.~Ghasemipour, R.~Gontijo-Lopes, B.~K.~Aryabhata, R.~Salimans,J.~Ho, D.~J.~Fleet, and M.~Norouzi,``Photorealistic text-to-image diffusion models with deep languageunderstanding,''in \textit{Advances in NeurIPS}, vol.~35, 2022.

\bibitem{ramesh2022dalle2}A.~Ramesh, P.~Dhariwal, A.~Nichol, C.~Chu, and M.~Chen,``Hierarchical text-conditional image generation with CLIP latents,''\textit{arXiv preprint arXiv:2204.06125}, 2022.

\bibitem{chefer2023attend}H.~Chefer, Y.~Alaluf, Y.~Vinker, L.~Wolf, and D.~Cohen-Or,``Attend-and-excite: Attention-based semantic guidance fortext-to-image diffusion models,''\textit{ACM Transactions on Graphics}, vol.~42, no.~4, 2023.

\bibitem{bau2019seeing}D.~Bau, J.~Zhu, H.~Strobelt, B.~Zhou, J.~B.~Tenenbaum,W.~T.~Freeman, and A.~Torralba,``Seeing what a GAN cannot generate,''in \textit{Proceedings of IEEE/CVF ICCV}, 2019.\end{thebibliography}

\end{document}